\section{The 1D diffusion equation}\label{sec:1d}

The first instantiation of the framework is the one-dimensional diffusion
equation, $d = 1$, with a single scalar component; we drop the component index
throughout this section. It serves as a testbed: every operator the framework
needs is available in closed form, so the only approximation in a run is the
operator splitting itself.

\subsection{Governing equation and clock}\label{ssec:1d-clock}

The dimensional equation is
\begin{equation}\label{eq:1d-dimensional}
  \frac{\partial u^*}{\partial t^*} = D \frac{\partial^2 u^*}{\partial (x^*)^2},
\end{equation}
so that, in the notation of \eqref{eq:fw-general-pde},
\begin{equation}\label{eq:1d-R}
  R^*(u^*) = D \frac{\partial^2 u^*}{\partial (x^*)^2}.
\end{equation}
Under the rescaling \eqref{eq:fw-nondim} this term becomes
\begin{equation}\label{eq:1d-R-rescaled}
  D \frac{\partial^2 (u U)}{\partial (x^*)^2}
  = D U \frac{\partial}{\partial x}\!\left( \frac{\partial u}{\partial x} \frac{1}{L} \right) \frac{1}{L}
  = \frac{D U}{L^2} \frac{\partial^2 u}{\partial x^2},
\end{equation}
having used $\mathrm{d}x/\mathrm{d}x^* = 1/L$. The clock is then chosen to make
the physical term of the lemma \eqref{eq:fw-lemma} carry unit coefficient,
\begin{equation}\label{eq:1d-clock}
  \frac{\tau}{U} \cdot \frac{D U}{L^2} = 1
  \quad\Longrightarrow\quad
  \tau = \frac{L^2}{D},
\end{equation}
which is the diffusive time scale of the current length scale, as expected.
In log form, as required by \eqref{eq:fw-ln-dtstar},
\begin{equation}\label{eq:1d-ln-clock}
  \ln \tau_k = 2 \ln L_k - \ln D .
\end{equation}
Substituting \eqref{eq:1d-clock} into the rescaling lemma
\eqref{eq:fw-lemma} gives the equation the algorithm integrates,
\begin{equation}\label{eq:1d-transformed}
  \frac{\partial u}{\partial t}
  = \sigma_L\, x \frac{\partial u}{\partial x}
  - \sigma_U\, u
  + \frac{\partial^2 u}{\partial x^2}.
\end{equation}

\subsection{The physics step}\label{ssec:1d-physics-step}

Because the diffusion term in \eqref{eq:1d-transformed} has unit coefficient,
the physics operator \eqref{eq:fw-physics-flow},
\begin{equation}\label{eq:1d-physics-flow}
  \Phi_{\Delta t}^{(\mathrm{P})}(u_k)
  = u_k + \int_{t_k}^{t_k + \Delta t} \frac{\partial^2 u}{\partial x^2}\, \mathrm{d}t,
\end{equation}
has the closed-form spectral solution
\begin{equation}\label{eq:1d-propagator}
  \hat u_j^{(\mathrm{P})} = \hat u_j\, e^{-k_j^2 \Delta t},
  \qquad k_j = \frac{2\pi j}{H}, \quad j = 0, \dots, \left\lfloor N/2 \right\rfloor,
\end{equation}
where $\hat u_j = \mathrm{RFFT}(u)_j$. This is an exact one-shot spectral
propagator, not an approximate time-marching scheme, so no sub-stepping or
implicit solve is needed, and the accuracy requirement on \textsc{Evolve}
(\S\ref{ssec:alg-substeps}) is met trivially at any order. Inverse-transforming
gives the advanced field $u^{(\mathrm{P})}(x_n)$.

\subsection{Moment constraints}\label{ssec:1d-moments}

The constraint functionals of \S\ref{ssec:fw-constraints} are taken to be the
zeroth and second moments of the field,
\begin{equation}\label{eq:1d-functionals}
  F_1[u] = \int u\, \mathrm{d}x \equiv M_0,
  \qquad
  F_2[u] = \int x^2 u\, \mathrm{d}x \equiv M_2,
\end{equation}
so that $M_0$ measures the amount of the conserved quantity and $M_2$ its
spread. Writing $M_0^*, M_2^*$ for the moments of the field being scaled, and
substituting $y = x\, \tilde L$ in each integral, the scaled field carries
\begin{equation}\label{eq:1d-M0-homogeneity}
  \int \frac{1}{\tilde U} u\!\left( x\, \tilde L \right) \mathrm{d}x
  = \frac{1}{\tilde U \tilde L}\, M_0^*,
\end{equation}
\begin{equation}\label{eq:1d-M2-homogeneity}
  \int x^2\, \frac{1}{\tilde U} u\!\left( x\, \tilde L \right) \mathrm{d}x
  = \frac{1}{\tilde U \tilde L^{3}}\, M_2^*,
\end{equation}
so the homogeneity exponents of \eqref{eq:fw-homogeneity} are
\begin{equation}\label{eq:1d-exponents}
  (p_1, q_1) = (1, 1), \qquad (p_2, q_2) = (1, 3),
\end{equation}
with $p_1 q_2 - p_2 q_1 = 2 \neq 0$, so the constraint system
\eqref{eq:fw-constraint-system} is uniquely solvable: the two moments do
respond independently to amplitude and to length rescaling. Solving it,
\begin{equation}\label{eq:1d-Ltilde}
  \tilde L = \sqrt{\frac{M_0\, M_2^*}{M_0^*\, M_2}},
\end{equation}
\begin{equation}\label{eq:1d-Utilde}
  \tilde U = \frac{M_0^*}{M_0\, \tilde L}
  = \left(\frac{M_0^*}{M_0}\right)^{3/2} \sqrt{\frac{M_2}{M_2^*}} .
\end{equation}

\emph{Open.} An earlier form of \eqref{eq:1d-Ltilde} and \eqref{eq:1d-Utilde}
omitted the factor $M_0$, giving
$\tilde L = \sqrt{M_2^*/(M_0^* M_2)}$ and
$\tilde U = (M_0^*)^{3/2} M_0^{-1} \sqrt{M_2/M_2^*}$; the present
implementation uses that form. The two agree exactly when $M_0 = 1$, and
differ by a factor $M_0^{\pm 1/2}$ otherwise. The discrepancy traces to
whether the target second moment is taken raw, as in \eqref{eq:1d-functionals}
and \eqref{eq:alg-targets}, or normalized by $M_0$. Since diffusion conserves
$M_0$ and the initial condition is in practice normalized to unit mass, the
difference has not so far been visible in a run. Which convention is intended
should be settled, and the implementation brought into line with it.

\subsection{Spectral interpolation on the real transform}\label{ssec:1d-interpolation}

For $d = 1$ and a real-valued field, the generic interpolation formula
\eqref{eq:fw-spectral-interp} is realized on the non-negative frequency
components of a real Fast Fourier Transform ($\mathrm{RFFT}$):
\begin{equation}\label{eq:1d-interp}
  \begin{split}
    \Phi_{\Delta t}^{(\mathrm{S})}(u_k)(x) = \frac{1}{\tilde{U}_{k+1}}\biggl\{ &\hat{u}_0^{(k)}
      + 2 \sum_{j=1}^{N/2-1} \biggl[ \operatorname{Re}\bigl(\hat{u}_j^{(k)}\bigr) \cos\left(\frac{2\pi j \tilde{L}_{k+1} x}{H}\right) \\
      &\qquad - \operatorname{Im}\bigl(\hat{u}_j^{(k)}\bigr) \sin\left(\frac{2\pi j \tilde{L}_{k+1} x}{H}\right) \biggr] \\
      &+ \hat{u}_{N/2}^{(k)} \cos\left(\frac{\pi N \tilde{L}_{k+1} x}{H}\right) \biggr\},
  \end{split}
\end{equation}
where $\hat{u}_j^{(k)} = \frac{1}{N} \mathrm{RFFT}(u_k)_j$ for
$j = 0, 1, \dots, N/2$, and the Nyquist term at $j = N/2$ is included only if
$N$ is even. The factor of $2$ accounts for the Hermitian conjugate modes
omitted by the real transform.

\subsection{Boundary treatment}\label{ssec:1d-boundary}

The substep \textsc{ApplyBC} of \S\ref{ssec:alg-substeps} imposes a homogeneous
Dirichlet condition by clamping the two edge grid points of the centred domain
\eqref{eq:alg-grid} after each completed step,
\begin{equation}\label{eq:1d-clamp}
  u_{k+1}(x_0) \gets 0, \qquad u_{k+1}(x_{N-1}) \gets 0 .
\end{equation}
On the centred grid $x \in [-H/2, H/2)$ these are two distinct grid points, not
periodic images of one another.

\emph{Open.} This sits in tension with the rest of the discretization: the
spectral propagator \eqref{eq:1d-propagator} and the spectral interpolation
\eqref{eq:1d-interp} both presume a periodic field, whereas
\eqref{eq:1d-clamp} imposes a Dirichlet condition. The two are consistent only
so long as $u$ is negligible near the domain edges, which requires the initial
condition to be localized well inside the domain and the run to be stopped
before the profile spreads to the boundary. The scaling step works against
this --- it is what keeps the profile from spreading in the self-similar frame
--- but the condition is an assumption of the scheme rather than a consequence
of it, and is not currently monitored.
